% Generated draft: requires mathematical and source review.
\documentclass[11pt]{article}
\usepackage[margin=25mm]{geometry}
\usepackage{fontspec}
\setmainfont{Latin Modern Roman}
\usepackage{amsmath,amssymb,mathtools,unicode-math}
\setmathfont{Latin Modern Math}
\usepackage{enumitem,needspace,xurl,hyperref}
\setlength{\parindent}{0pt}
\setlength{\parskip}{5pt}
\setlength{\emergencystretch}{3em}
\newcommand{\N}{\mathbb{N}}
\newcommand{\Z}{\mathbb{Z}}
\newcommand{\Q}{\mathbb{Q}}
\newcommand{\R}{\mathbb{R}}
\newcommand{\C}{\mathbb{C}}
\newcommand{\F}{\mathbb{F}}
\newcommand{\A}{\mathbb{A}}
\newcommand{\PP}{\mathbb{P}}
\newcommand{\E}{\mathbb{E}}
\newcommand{\eps}{\varepsilon}
\newcommand{\dd}{\,\mathrm{d}}
\begin{document}
\section*{UnsolvedMath ID: 700943}
\textbf{Source problem code:} catalogue entry
\subsection{Definitions and mathematical statement}
Let \((G,\tau)\) be a topological group, with identity element \(e\). An ultrafilter on the underlying set of \(G\) is a family \(p\subseteq\mathcal P(G)\) such that \(\varnothing\notin p\), \(G\in p\), \(p\) is closed under finite intersections and passage to supersets, and for every \(A\subseteq G\), exactly one of \(A\) and \(G\setminus A\) belongs to \(p\). It converges to \(e\) if every \(\tau\)-neighborhood of \(e\) belongs to \(p\).

For ultrafilters \(p,q\) on \(G\), use the standard product determined by
\[
 A\in p\cdot q
 \quad\Longleftrightarrow\quad
 \{x\in G:x^{-1}A\in q\}\in p,
 \qquad A\subseteq G,
\]
where \(x^{-1}A=\{y\in G:xy\in A\}\). Under the usual interpretation, write
\[
 \tau(G)=\{p:p\text{ is an ultrafilter on }G\text{ converging to }e\},
\]
with this multiplication. The topology here is the given topology \(\tau\); the ultrafilters are on the full power set of the underlying group.

\textbf{Question (E. G. Zelenyuk, 1995).} For every topological group \((G,\tau)\) such that \(\tau(G)\) is finite, must
\[
 \forall p\in\tau(G),\qquad p\cdot p=p
\]
hold? Thus the requested conclusion is that every element of \(\tau(G)\) is idempotent; commutativity is not requested.

\textbf{Qualifications requiring review.} The supplied record refers to Protasov's paper for the ultrafilter semigroup construction without giving its definitions. The displayed multiplication is the standard interpretation and should be checked against that reference. The record does not specify whether the principal ultrafilter at \(e\) is included. Including or excluding that ultrafilter does not change finiteness or the assertion that every element is idempotent, since it is the identity for ultrafilter multiplication. No separation axiom for the topological group is explicitly supplied; none is added here.

\textbf{Supplied comment of 2005.} The record reports that the assertion was proved when the underlying group \(G\) is countable, citing E. G. Zelenyuk (2000). This is a supplied historical qualification, not an independently verified result.
\subsection{Short English statement}
E. G. Zelenyuk asks whether a finite semigroup of ultrafilters converging to the identity of a topological group must consist entirely of idempotents. The supplied 2005 comment reports a positive answer for countable groups. The precise ultrafilter convention and any intended separation assumptions require source review.
\subsection{Sources}
\begin{itemize}
\item Local UnsolvedMath registry: permanent ID \(700943\), exact source problem code \texttt{catalogue entry}, title \emph{Kourovka Notebook Problem 13.25}. The supplied registry status is \texttt{open}; this label has not been independently verified.
\item \emph{The Kourovka Notebook (Unsolved Problems in Group Theory)}, problem 13.25, proposed by E. G. Zelenyuk in the 13th issue (1995). The supplied background identifies the consulted edition as the 21st issue edition (2026), arXiv:1401.0300v48. The supplied collection description names E. I. Khukhro and V. D. Mazurov as editors.
\item I. V. Protasov, \emph{Siberian Math. J.}, \textbf{34}, no. 5 (1993), 938--951. Cited in the supplied statement for the ultrafilter semigroup construction; no article title or URL is supplied.
\item E. G. Zelenyuk, \emph{Mat. Studii}, \textbf{14} (2000), 121--140. Cited by the supplied comment of 2005 for the countable-group case; no article title or URL is supplied.
\end{itemize}
\end{document}
